\documentclass[11pt,a4paper]{article}
\usepackage[utf8]{inputenc} % Not needed with XeLaTeX
\usepackage{fontspec}
\usepackage[margin=0.75in]{geometry}
\usepackage{xcolor}
\definecolor{darknavy}{RGB}{0,0,100}
\usepackage[colorlinks=true, urlcolor=darknavy, linkcolor=darknavy]{hyperref}
\usepackage{enumitem}
% \usepackage{fontawesome}
\usepackage{fontawesome5}
\usepackage{titlesec}
\usepackage{academicons}

% % --- Font: Source Sans Pro via Type1 (works with pdfLaTeX)
% \usepackage[default,scale=1.0]{sourcesanspro} % body + sans by default
% \renewcommand{\familydefault}{\sfdefault}      % make sans the default family
% \usepackage[scaled]{helvet}
\usepackage[default,scale=0.95]{opensans}

% (optional) math that pairs reasonably with sans body
\usepackage{newtxsf} % sans math; or use newtxtext/newtxmath for serif math

\titleformat{\section}{\Large\bfseries}{\thesection}{1em}{}[\titlerule]
\titlespacing{\section}{0pt}{12pt}{6pt}

\usepackage{comicneue}

\begin{document}

\begin{center}
  % {\Huge\textbf{Pavel Popov}}\\[0.5em]
  % {\comicneue\Huge\textbf{Pavel Popov}}\\[0.5em]
  {\fontspec{JetBrains Mono}\Huge Pavel Popov}\\[0.5em]
  {\fontspec{JetBrains Mono}\footnotesize +1 (470) 920-1876 \quad | \quad Atlanta, GA 30318} \quad | \quad {paavali.popov@gmail.com}\\
  {\small \href{https://github.com/paavalipopov}{Github: paavalipopov} \quad | \quad
  % \href{https://www.linkedin.com/in/pavel-popov-83b69a123/}{\faLinkedin\ Pavel Popov} \quad
  \href{https://www.linkedin.com/in/pavel-popov-83b69a123/}{LinkedIn: pavel-popov-83b69a123} \quad | \quad
  \href{https://scholar.google.com/citations?user=4alUD5UAAAAJ}{Google scholar: 4alUD5UAAAAJ}}
\end{center}

\section*{Skills}
\begin{itemize}
[leftmargin=*,noitemsep]
\item \textbf{Visualization \& Data Exploration:} Data visualization (Plotly, Matplotlib, Seaborn), interactive visualization, latent-space \& representation analysis, Figma.
\item \textbf{Machine Learning:} Deep learning architecture design, representation learning, explainable/interpretable AI, time series analysis, causal inference, distributed data parallel training.
\item \textbf{Data \& Infrastructure:} Large-scale multimodal imaging \& time-series datasets, dataset benchmarking \& management, feature engineering, SQL, mongoDB, Spark.
\item \textbf{Programming \& Tools:} Python, PyTorch, Scikit-learn, Pandas, MATLAB, Git, Docker, HPC, SLURM.
\end{itemize}

\section*{Open Source Projects}
\noindent\textbf{\href{https://pypi.org/project/ml4fmri/}{ml4fmri}: Python toolkit for benchmarking \& exploring DL models on large-scale time-series data}\\
[-2em]
\begin{itemize}[leftmargin=*,noitemsep]
    \item Unified 13 DL models for multivariate time series analysis under a standardized training/evaluation API, giving ML researchers a single interface to navigate and compare models across datasets.
    \item Automated an end-to-end benchmarking pipeline (ingestion, training, evaluation) for reproducible research on arbitrary input time series.
\end{itemize}

\noindent\textbf{\href{https://pypi.org/project/brainbow/}{brainbow} --- Visualization tool for multi-dimensional neuroimaging data}\\
[-2em]
\begin{itemize}[leftmargin=*,noitemsep]
    \item Built a command-line visualization tool for rendering and exploring 3D/4D neuroimaging data formats.
\end{itemize}

\section*{Experience}
\noindent{\textbf{Center for Translational Research in Neuroimaging and Data Science}} \hfill January 2022 --- Now\\
Graduate Research Assistant \hfill \textit{Atlanta, USA}\\
[-2em]
\begin{itemize}[leftmargin=*,noitemsep]
    \item Built researcher-facing ML tooling and reproducible benchmarking frameworks for large-scale multimodal imaging \& time-series datasets (4D fMRI --- volumetric image/video data), standardizing how researchers train, evaluate, and compare deep learning models.
    \item Developed a robust DL classifier for MRI image diagnostics with SOTA performance and 10x speedup over existing benchmarks.
    \item Designed interpretable, latent-space analysis methods to surface semantic representations and novel biomarkers from deep models on multivariate imaging \& time-series data.
    \item Built and optimized distributed training and deployment pipelines (HPC, SLURM, distributed data parallel) for DL on large imaging datasets.
    \item Collaborated with cross-functional AI and clinical teams to translate research workflows into shared tooling; presented findings at internal meetings and international conferences.
\end{itemize}

\noindent\textbf{Institute of Radioengineering and Electronics} \hfill May 2017 --- September 2020\\
Research Assistant \hfill \textit{Moscow, Russia}\\
[-2em]
\begin{itemize}[leftmargin=*,noitemsep]
    \item Modeled complex signal propagation in novel hardware architectures based on wave logic.
    \item Developed mathematical models for control of high-frequency signal dynamics, providing theoretical validation for prototype hardware controllers.
    \item Ran large-scale numerical simulations to validate theoretical models.
\end{itemize}

\section*{Education}
\noindent\textbf{Georgia State University} \hfill June 2027\\
PhD in Computer Science \hfill
% Machine Learning \& Neuroimaging,
GPA 3.89\\[0.3em]
\textbf{Moscow Institute of Physics and Technology} \hfill June 2020\\
MS in Applied Mathematics and Physics \hfill
% Spintronics,
GPA 3.92\\[0.3em]
\textbf{Moscow Institute of Physics and Technology} \hfill June 2018\\
BS in Applied Mathematics and Physics \hfill
% Solid State Physics,
GPA 3.81

\end{document}
